% Archived 2026-10 when langgraph_intro.tex was trimmed for the one-hour Seminar.
% Part 1: duplicate frames dropped. Part 2: commented-out blocks from the old file.
%%%%%%%%%%%%%%%%%%%%%%%%%%%%%%%%%%%%%%%%%%%%%%%%%%%%%%%%%%%%%%%%%%%%%%%%%%%%%%%%
\begin{frame}[fragile]\frametitle{Introduction to LangGraph}
    \begin{itemize}
        \item LangGraph extends LangChain for complex, stateful workflows.
        \item Uses \textbf{State Graphs} with nodes (tasks) and edges (flow logic).
        \item Enables loops, conditional branching, and persistent state memory.
        \item Ideal for multi-step, research-style assistants beyond chatbots.
        \item Each node handles a specific subtask (search, scrape, evaluate, report).
        \item Shared state retains context like URLs, content, scores, and facts.
        \item Simplifies orchestration: LangGraph manages task sequencing.
        \item Focus shifts from coding logic to workflow architecture.
    \end{itemize}
\end{frame}
%%%%%%%%%%%%%%%%%%%%%%%%%%%%%%%%%%%%%%%%%%%%%%%%%%%%%%%%%%%%%%%%%%%%%%%%%%%%%%%%
\begin{frame}[fragile]\frametitle{LangGraph in Practice}
    \begin{itemize}
        \item Nodes: search, scrape, evaluate trust, extract facts, generate report.
        \item State persists across nodes, similar to long-term memory.
        \item Conditional edges allow adaptive decision-making in workflows.
        \item Loops enable iterative improvement of research quality.
        \item Tool nodes integrate external data sources (e.g., DuckDuckGo).
        \item Memory nodes accumulate knowledge and generate comprehensive reports.
        \item Demonstrates scalable, intelligent workflow automation for enterprises.
    \end{itemize}
\end{frame}
%%%%%%%%%%%%%%%%%%%%%%%%%%%%%%%%%%%%%%%%%%%%%%%%%%%%%%%%%%%%%%%%%%%%%%%%%%%%%%%%
\begin{frame}[fragile]\frametitle{How is LangGraph Different?}


      \begin{itemize}
        \item Gives you more control over how your AI makes decisions
        \item Allows your AI to revisit previous steps if needed
        \item Makes it easier to add human oversight at specific points
        \item Integrates smoothly with other LangChain tools
      \end{itemize}


\begin{center}
\includegraphics[width=\linewidth,keepaspectratio]{langgraph15}
\end{center}

\end{frame}
%%%%%%%%%%%%%%%%%%%%%%%%%%%%%%%%%%%%%%%%%%%%%%%%%%%%%%%%%%%%%%%%%%%%%%%%%%%%%%%%
\begin{frame}[fragile]\frametitle{Conditional Edges in LangGraph}
      \begin{itemize}
        \item Conditional edges enable branching logic within the graph.
        \item After a node executes, a function decides the next path.
        \item Typically implemented using a \texttt{should\_continue} function.
        \item The function returns a key such as “continue” or “end”.
        \item The returned key maps to the next node in the graph.
        \item “continue” → proceed to the next action node.
        \item “end” → trigger the built-in end node and stop execution.
        \item Enables dynamic control flow based on runtime data.
      \end{itemize}
\end{frame}
%%%%%%%%%%%%%%%%%%%%%%%%%%%%%%%%%%%%%%%%%%%%%%%%%%%%%%%%%%%%%%%%%%%%%%%%%%%%%%%%
\begin{frame}[fragile]\frametitle{State Management in LangGraph}
      \begin{itemize}
        \item State: Shared memory object accessible across graph nodes.
        \item Schema defines all possible fields before runtime.
        \item No new fields can be added at runtime; the schema is fixed.
        \item Each node can push updates to the state dynamically.
        \item Updates modify existing fields via reducer functions.
        \item Default update behavior: overwrite existing field values.
        \item Can also accumulate data (e.g., append to lists).
        \item Encourages structured and predictable data flow across nodes.
      \end{itemize}
\end{frame}
%%%%%%%%%%%%%%%%%%%%%%%%%%%%%%%%%%%%%%%%%%%%%%%%%%%%%%%%%%%%%%%%%%%%%%%%%%%%%%%%
\begin{frame}[fragile]\frametitle{LangGraph Key Features}
      \begin{itemize}
        \item \textbf{Looping and Branching}: Conditional statements and loop structures
        \item \textbf{State Persistence}: Automatic save/restore, pause and resume
        \item \textbf{Human-in-the-Loop}: Insert human review, state editing capabilities
        \item \textbf{Streaming Processing}: Real-time feedback on execution status
        \item \textbf{LangChain Integration}: Reuses existing components, LCEL support
        \item Provides controlled flexibility unlike pure React agents
        \item Production-ready with debugging and observability
      \end{itemize}
\end{frame}
%%%%%%%%%%%%%%%%%%%%%%%%%%%%%%%%%%%%%%%%%%%%%%%%%%%%%%%%%%%%%%%%%%%%%%%%%%%%%%%%
\begin{frame}[fragile]\frametitle{Step 3: Define Routing Logic}

      \begin{itemize}
        \item Routing function decides next node based on state
        \item Returns the name of the target node
        \item Enables dynamic workflow execution
      \end{itemize}
	  
      \begin{lstlisting}
def route_email(state: EmailState) -> str:
    """Conditional routing based on classification"""
    if state["is_spam"]:
        return "spam"
    return "not_spam"
      \end{lstlisting}
     
\end{frame}
%%%%%%%%%%%%%%%%%%%%%%%%%%%%%%%%%%%%%%%%%%%%%%%%%%%%%%%%%%%%%%%%%%%%%%%%%%%%%%%%
\begin{frame}[fragile]\frametitle{Summary}
      \begin{itemize}
        \item LangGraph enables controlled, flexible AI workflows through graph structures
        \item Solves limitations of rigid chains and unpredictable pure agents
        \item Key components: Nodes (logic), Edges (flow), State (context)
        \item Ideal for workflow automation with intelligent decision-making
        \item Production-ready with state persistence, human-in-loop, streaming
        \item Start with simple workflows, expand to complex multi-agent systems
        \item The future of reliable, autonomous AI applications
      \end{itemize}
\end{frame}
%%%%%%%%%%%%%%%%%%%%%%%%%%%%%%%%%%%%%%%%%%%%%%%%%%%%%%%%%%%%%%%%%%%%%%%%%%%%%%%%
% ---- commented-out blocks from the old langgraph_intro.tex ----
% %%%%%%%%%%%%%%%%%%%%%%%%%%%%%%%%%%%%%%%%%%%%%%%%%%%%%%%%%%%
% \begin{frame}[fragile]\frametitle{Level 1-2: Code and Single LLM Calls}
      % \begin{itemize}
        % \item \textbf{Level 1: Pure Code}:
        % \begin{itemize}
            % \item Example: \lstinline|if temperature > 30: return "It's hot"| 
            % \item No autonomy, 100\% deterministic, predictable
            % \item Problem: Must anticipate every scenario
            % \item Use case: Simple rule-based logic
        % \end{itemize}
        % \item \textbf{Level 2 - Single LLM Call}:
        % \begin{itemize}
            % \item Example: "Summarize this article" → LLM → Summary
            % \item One input one output. One atomic task, limited reasoning
            % \item Problem: Trying to do everything in one shot. Cannot break down complex tasks
            % \item Use case: Text generation, classification, extraction
        % \end{itemize}
        % \item \textbf{Limitation}: Both lack multi-step reasoning capability
      % \end{itemize}
% \end{frame}
% %%%%%%%%%%%%%%%%%%%%%%%%%%%%%%%%%%%%%%%%%%%%%%%%%%%%%%%%%%%
% \begin{frame}[fragile]\frametitle{Level 3: Chains: Sequential Reasoning}
      % \begin{itemize}
        % \item \textbf{What}: Unidirectional sequence of LLM operations, so multiple experts one after another. Breaking down task is possible now.
        % \item \textbf{Example: Customer Support}:
        % \begin{itemize}
            % \item Step 1: Classify intent (complaint/question/request)
            % \item Step 2: Extract key information (order ID, issue)
            % \item Step 3: Generate appropriate response
        % \end{itemize}
        % \item \textbf{Flow}: A → B → C → End (no branching, no loops)
        % \item \textbf{Advantages}: Predictable, reliable, easy to debug
        % \item \textbf{Limitations}: 
        % \begin{itemize}
            % \item Cannot adapt to different scenarios
            % \item No conditional logic or branching
            % \item Fixed execution path regardless of context
        % \end{itemize}
        % \item Like an assembly line, efficient but inflexible
      % \end{itemize}
% \end{frame}
% %%%%%%%%%%%%%%%%%%%%%%%%%%%%%%%%%%%%%%%%%%%%%%%%%%%%%%%%%%%
% \begin{frame}[fragile]\frametitle{Level 4: Routers: Conditional Branching}
      % \begin{itemize}
        % \item \textbf{What}: Chains with decision points, but still unidirectional
        % \item \textbf{Example: Document Processing}:
        % \begin{itemize}
            % \item Analyze document type, LLM/ML, decision making (invoice/contract/email)
            % \item IF invoice → Extract line items → Calculate total
            % \item IF contract → Extract parties → Identify key terms
            % \item IF email → Classify urgency → Route to handler
        % \end{itemize}
        % \item \textbf{Flow}: Decision node routes to specialized chains
        % \item \textbf{Advantages}: Handles different scenarios, specialized processing
        % \item \textbf{Limitations}:
        % \begin{itemize}
            % \item Still no loops or backward flow
            % \item Cannot learn from mistake, cannot revisit back
            % \item One-way path through the system
        % \end{itemize}
        % \item Like a train switching tracks - more flexible but still one direction
      % \end{itemize}
% \end{frame}
% %%%%%%%%%%%%%%%%%%%%%%%%%%%%%%%%%%%%%%%%%%%%%%%%%%%%%%%%%%%
% \begin{frame}[fragile]\frametitle{Level 5: State: Flow with Loops}
      % \begin{itemize}
        % \item \textbf{What}: State machines are Routers who can branch and loop.
        % \item \textbf{Example: Content Creator}:
        % \begin{itemize}
            % \item Given a theme to write on
			% \item HEAD can sending it for LinkedIn article writing, 
			% \item Once done, we can have Human-approval, if not approved go back to the HEAD. Revise, iterate till approved.
			% \item Once approved send to other forms such as Blog post, Twitter announcement post etc.
        % \end{itemize}
        % \item \textbf{Key Features}: 
        % \begin{itemize}
            % \item Can have human-in-loop for approval
			% \item Multi-Agent systems
			% \item Advanced Memory Management
			% \item Can go back in history and explore better/alternate paths
        % \end{itemize}		
        % \item \textbf{Challenge}: Flow has to be prescribed, not full autonomy.
        % \item \textbf{This is LangGraph}
      % \end{itemize}
% \end{frame}
% %%%%%%%%%%%%%%%%%%%%%%%%%%%%%%%%%%%%%%%%%%%%%%%%%%%%%%%%%%%%%%%%%%%%%%%%%%%%%%%%%%
% \begin{frame}[fragile]\frametitle{Level 5: State: Flow with Loops}
% \begin{center}
% \includegraphics[width=0.9\linewidth,keepaspectratio]{langgraph24}
% {\tiny (Ref: LangGraph Crash Course - Harish Neel)}
% \end{center}	  
% \end{frame}
% %%%%%%%%%%%%%%%%%%%%%%%%%%%%%%%%%%%%%%%%%%%%%%%%%%%%%%%%%%%
% \begin{frame}[fragile]\frametitle{Level 6: Agents: Autonomous with Loops}
      % \begin{itemize}
        % \item \textbf{What}: State machines with independent thinking, can loop and branch
        % \item \textbf{Example - Research Assistant}:
        % \begin{itemize}
            % \item User: "Analyze SpaceX's latest launch success rate"
            % \item Agent thinks: "I need current data and historical context"
            % \item Action 1: Search web for latest SpaceX launches
            % \item Observation: Found data but it's incomplete
            % \item Action 2: Search for historical launch records
            % \item Observation: Now have full picture
            % \item Final answer: Synthesize findings into analysis
        % \end{itemize}
        % \item \textbf{Key Features}: Think → Act → Observe → Decide → Repeat
        % \item \textbf{Challenge}: Can loop infinitely if not controlled
      % \end{itemize}
% \end{frame}
% %%%%%%%%%%%%%%%%%%%%%%%%%%%%%%%%%%%%%%%%%%%%%%%%%%%%%%%%%%%%%%%%%%%%%%%%%%%%%%%%%%
% \begin{frame}[fragile]\frametitle{Levels of Autonomy in LLM Applications}
% \begin{center}
% \includegraphics[width=0.8\linewidth,keepaspectratio]{langgraph5}
% {\tiny (Ref: LangGraph Crash Course - Harish Neel)}
% \end{center}	  
% LLM (What AI decides) vs Code (What we prescribe/program)
% \end{frame}
% # def classify_email(state: EmailState) -> EmailState:
% #     # WRONG - Mutating incoming state
% #     state.is_spam = True  # DON'T DO THIS
% #     state.messages.append(...)  # DON'T DO THIS
% #     return state
      % \end{lstlisting}
